% ==========================================================================
% GENERATED by scripts/split_paper.py. Do not edit this file.
%
% Source: imported_papers/grokking/iclr2027_conference_v2.tex
% Part 9: Controls on credit equalization
%
% Edits belong upstream, in the published paper. To change what the
% chapter includes, change imports/grokking.toml and re-run
% the splitter.
% ==========================================================================
\section{Controls on credit equalization}
\label{app:credit_controls}
The main experiment changes the credit magnitudes and observes a change in the outgoing weights. This section tests the alternative explanations: that the effect comes from the incoming weight norms rather than from the credit, that it depends on which neurons receive the credit, that Adam or the hyperparameters produce it, and that the concentrated circuit is only a transient state.

\subsection{Constraining incoming weights}
\label{app:named}
Full credit equalization gives every neuron the same gradient norm, and Nero, which also distributes the circuit, additionally constrains the weight norms. Could equal incoming weight norms alone explain the redistribution? We test two constraints on the weights that leave the gradients unchanged. The first renormalizes each neuron's incoming weights to a fixed norm after every update. Circuit participation is $0.661\,(0.015)$, compared with $0.595\,(0.023)$ for the reference, and deletion tolerance decreases (Table~\ref{tab:named}, upper block). The second is weight normalization \citep{salimans2016weightnorm}, which parameterizes each incoming weight vector as a direction times a learned scale; with weight decay on the scales it groks to test accuracy $0.993$ with participation $0.466\,(0.026)$. Neither constraint approaches the participation of $0.99$ reached by credit equalization. Equal incoming weight norms do not reproduce the redistribution.

\subsection{Assigning credit independently of neuron identity}
\label{app:imposed}
In the first hidden layer, a neuron's incoming-weight gradient is the input multiplied by the backpropagated error, and that error depends linearly on the neuron's outgoing weights. Credit and outgoing norms are therefore coupled at every step, and their agreement could reflect this coupling rather than an effect of the credit on the circuit. We break the coupling by assigning credit magnitudes that carry no information about the neurons.

A \emph{credit profile} is a vector of $n$ nonnegative target magnitudes with a chosen participation ratio. At every step we replace the incoming-gradient norms by this profile, scaled so that the total norm is preserved, and keep every gradient's direction. A random permutation, drawn once per run, decides which neuron receives which magnitude, so the assignment is independent of the neurons' roles. Circuit participation follows the imposed value, from $0.637$ at profile participation $0.60$ to $0.990$ at $1.00$, deletion tolerance increases with it, and test accuracy stays above $0.96$ (Table~\ref{tab:named}, lower block). An allocation chosen independently of the neurons' current roles therefore sets the final concentration.

A second control keeps the reference's credit magnitudes at each step but randomly reassigns them to the neurons at every step. Each permutation preserves that step's participation and changes only which neurons receive the largest signals. Final circuit participation is $0.992$, compared with $0.613$ for the reference. Measured over the last $800$ epochs, credit participation is $0.979$ rather than the reference's $0.604$: once the assignment no longer favors the same neurons, the trained network itself develops more uniform gradients. This suggests that concentration requires the credit to stay concentrated on the same neurons over time.

The effect persists across the tested initialization scales. Across scales from $0.5$ to $4$, reference participation is $0.586$ to $0.613$ and full equalization gives $0.990$ to $0.998$, with a mean absolute difference of $0.008$ between credit and circuit participation. The effect is also local to the layer in which credit is equalized: equalizing only the second hidden layer gives participation $1.000$ there while the first layer stays at $0.593\,(0.004)$ with $86$ low-norm neurons, and equalizing both layers distributes both, at test accuracy of at least $0.98$.

\begin{table}[!ht]
\caption{Controls on the weights and on the credit profile ($p=31$, three seeds; standard deviations in parentheses). Upper block: constraining the incoming weight norms without rescaling the gradients leaves the circuit concentrated. Lower block: credit profiles of a chosen participation, assigned to neurons by a random permutation, produce circuit participation within $0.018$ of the imposed value on average; the final row reassigns the reference's credit magnitudes to random neurons at every step. Each block has its own reference run. Random del and norm-ranked del are deletion tolerances (Appendix~\ref{app:deletion}).}
\label{tab:named}
\begin{center}
\small
\setlength{\tabcolsep}{4pt}
\begin{tabular}{lccccc}
\toprule
Condition & Test & Credit PR & Circuit PR & Random del & Norm-ranked del \\
\midrule
\multicolumn{6}{l}{\emph{Constraining the incoming weights}} \\
Reference & $0.998$ & $0.586$ & $0.595\,(0.023)$ & $0.672$ & $0.311$ \\
Fixed incoming norm & $0.954$ & $0.611$ & $0.661\,(0.015)$ & $0.607$ & $0.256$ \\
\midrule
\multicolumn{6}{l}{\emph{Imposing a credit profile}} \\
Reference & $0.997$ & $0.604$ & $0.613\,(0.016)$ & $0.673$ & $0.298$ \\
Imposed $0.60$ & $0.980$ & $0.600$ & $0.637\,(0.003)$ & $0.607$ & $0.240$ \\
Imposed $0.70$ & $0.981$ & $0.700$ & $0.727\,(0.001)$ & $0.663$ & $0.319$ \\
Imposed $0.80$ & $0.972$ & $0.800$ & $0.815\,(0.005)$ & $0.693$ & $0.420$ \\
Imposed $0.90$ & $0.969$ & $0.900$ & $0.902\,(0.005)$ & $0.723$ & $0.533$ \\
Imposed $1.00$ & $0.975$ & $1.000$ & $0.990\,(0.003)$ & $0.775$ & $0.599$ \\
Reshuffled each step & $0.974$ & $0.979$ & $0.992\,(0.002)$ & $0.775$ & $0.601$ \\
\bottomrule
\end{tabular}
\end{center}
\end{table}

\subsection{Optimizer and hyperparameter controls}
\label{app:baselines}
Adam normalizes each parameter separately, so one might suspect that Adam either creates the redistribution or removes it. Neither is the case. With SGD with momentum in place of Adam, circuit participation is $0.53$ at $\beta = 0$ and $0.99$ at $\beta = 1$, with deletion tolerance $0.60$ and $0.76$ (Table~\ref{tab:baselines}). With Lion in place of Adam, full credit equalization gives participation $0.984\,(0.003)$ at test accuracy $0.974\,(0.010)$, whereas Lion alone gives about $0.5$. In the Adam runs, participation measured from Adam's effective step $m/\sqrt{v}$ is $0.66$ at $\beta = 0$ and $0.99$ at $\beta = 1$, so the difference between neurons survives Adam's transformation.

Coupled weight decay is added to the rescaled gradient, so equalization also changes the size of the decay term relative to each neuron's gradient, and Appendix~\ref{app:decay} shows that under Adam this accounts for the redistribution. Our Lion harness also adds the decay before the update. For SGD with coupled decay, rescaling a neuron's gradient by $s_j$ equals a per-neuron learning rate $s_j$ with decay $\lambda/s_j$, whose fixed points are those of the decay control, so both optimizers follow the same per-neuron condition. To test the intervention without this pathway, we repeat the comparison with AdamW, which applies weight decay to the weights directly. Weight decay $1$ is the only value in $\{10^{-3}, 10^{-2}, 10^{-1}, 1\}$ at which the AdamW reference groks within $20{,}000$ epochs. At that value, over $80{,}000$ epochs and five seeds measured every $5000$ epochs, equalization gives higher outgoing participation at all $80$ paired seed checkpoints, with $0.96$ against $0.91$ at the end of training. Activation-weighted and ablation participation are higher under equalization at $69$ and $74$ of the $80$ checkpoints. Both functional measures fluctuate strongly in both conditions. Over the last $30{,}000$ epochs they range from $0.02$ to $0.77$ for the reference and from $0.13$ to $0.98$ under equalization, and at the end of training they average $0.41$ and $0.36$ for the reference and $0.82$ and $0.73$ under equalization. Test accuracy lies between $0.94$ and $1.00$ in both conditions. The low functional values of the reference come from $20$ to $50$ first-layer neurons that become active on every input and add a constant offset to all logits, which zero-ablation counts as large effects. Replacing these neurons by their mean activation changes test accuracy by at most $0.003$. Computed over the other neurons and without the constant part of each activation, activation-weighted participation is $0.88$ for the reference and $0.97$ under equalization, higher in all ten seeds. The AdamW reference is thus distributed in both kinds of measure, and equalization distributes it further.

Ordinary tuning of the reference does not distribute the circuit either. Across three learning rates and two weight-decay values, every setting groks with circuit participation between $0.530$ and $0.626$ (Table~\ref{tab:bpgrid}). Credit participation ranges from $0.523$ to $0.613$ and differs from circuit participation by $0.009$ on average.

Dropout is the natural comparison for deletion tolerance, since it trains the network under random deletion. It distributes the circuit as well, to participation $0.95$ at rate $0.1$ and $0.99$ at rate $0.5$, and its random-deletion tolerance, $0.79$ and $0.82$, is close to the $0.77$ of full credit equalization, which never sees a deleted neuron during training. Two unrelated interventions thus reach the same participation and nearly the same tolerance, which supports the reading of Appendix~\ref{app:suff} that, in this network, deletion tolerance follows the distribution of outgoing norms however that distribution is produced. The interventions differ in what else they change. Dropout modifies the objective and raises the margin and the parameter norm (Table~\ref{tab:baselines}), whereas credit equalization changes only the learning rule and leaves loss, margin and norm within a few percent of the reference. Intermediate equalization strengths lie between the endpoints: at $\beta = 0.5$, circuit participation is $0.98$ and test accuracy $0.985$; at $\beta = 0.75$, they are $0.99$ and $0.977$.

\begin{table}[!ht]
\caption{Optimizer and regularizer controls ($p=31$, three seeds). Upper block: circuit participation and random-deletion tolerance under other optimizers and under dropout; margin and norm are given where measured. Lower block: participation ratio of the measured learning signal in the Adam runs, from the raw gradient and from Adam's effective step $m/\sqrt{v}$.}
\label{tab:baselines}
\begin{center}
\small
\begin{tabular}{lcccc}
\toprule
Setting & Circuit PR & Random del & Margin & Norm \\
\midrule
\multicolumn{5}{l}{\emph{Circuit participation}} \\
SGD with momentum, $\beta{=}0$ & $0.53$ & $0.60$ & --- & --- \\
SGD with momentum, $\beta{=}1$ & $0.99$ & $0.76$ & --- & --- \\
Lion, $\beta{=}1$ & $0.984\,(0.003)$ & --- & --- & --- \\
Dropout $0.1$ & $0.95$ & $0.79$ & $5.3$ & $26.8$ \\
Dropout $0.5$ & $0.99$ & $0.82$ & $6.8$ & $30.2$ \\
Credit equalization, $\beta{=}1$ (Adam) & $0.99$ & $0.77$ & $5.2$ & $26.6$ \\
\midrule
\multicolumn{5}{l}{\emph{Participation of the measured learning signal (Adam runs)}} \\
Raw gradient, $\beta{=}0$ & $0.59$ & --- & --- & --- \\
Effective step $m/\sqrt{v}$, $\beta{=}0$ & $0.66$ & --- & --- & --- \\
Effective step $m/\sqrt{v}$, $\beta{=}1$ & $0.99$ & --- & --- & --- \\
\bottomrule
\end{tabular}
\end{center}
\end{table}

\begin{table}[!ht]
\caption{Learning-rate and weight-decay grid for the reference ($p=31$, three seeds; standard deviations in parentheses). All settings grok and retain concentrated circuits, with credit participation close to circuit participation.}
\label{tab:bpgrid}
\begin{center}
\small
\begin{tabular}{lcccc}
\toprule
Learning rate & Weight decay & Test & Credit PR & Circuit PR \\
\midrule
$3\times10^{-4}$ & $10^{-3}$        & $0.998$ & $0.613$ & $0.626\,(0.012)$ \\
$3\times10^{-4}$ & $3\times10^{-3}$ & $0.983$ & $0.560$ & $0.571\,(0.027)$ \\
$10^{-3}$        & $10^{-3}$        & $0.997$ & $0.604$ & $0.613\,(0.016)$ \\
$10^{-3}$        & $3\times10^{-3}$ & $0.993$ & $0.523$ & $0.530\,(0.008)$ \\
$3\times10^{-3}$ & $10^{-3}$        & $0.993$ & $0.589$ & $0.594\,(0.024)$ \\
$3\times10^{-3}$ & $3\times10^{-3}$ & $0.988$ & $0.546$ & $0.553\,(0.022)$ \\
\bottomrule
\end{tabular}
\end{center}
\end{table}

\subsection{Longer training and switching the learning rule}
\label{app:long}
Is the concentrated reference circuit a stable solution, or a stage on the way to a distributed one? We extend training from $20{,}000$ to $200{,}000$ epochs. Reference participation stabilizes by approximately $40{,}000$ epochs and ends at $0.605\,(0.012)$; full equalization ends at $0.990\,(0.003)$ (Table~\ref{tab:long}). Longer training does not distribute the reference circuit.

We then train for $20{,}000$ epochs with one rule and continue for $30{,}000$ epochs with the other, with same-rule continuations as controls. Switching from the reference to full equalization raises participation from $0.613\,(0.016)$ to $0.991\,(0.000)$ within $5000$ epochs, at test accuracy $0.997$; continuing with the reference leaves it at $0.610\,(0.012)$. The reverse switch lowers participation from $0.990\,(0.003)$ to $0.722\,(0.016)$ and is still decreasing when the run ends, while its control stays at $0.990\,(0.003)$. Task and data are unchanged in all four runs; only the credit assignment changes. A fixed rule thus preserves its circuit over the tested horizon, and changing the credit allocation reorganizes the outgoing weights of an already learned circuit in either direction. Figure~\ref{fig:steer}b shows a separate run of the same protocol at $100$-epoch resolution: after the switch to equalization, participation exceeds $0.95$ within $300$ epochs, and after the reverse switch it reaches $0.67$ after $30{,}000$ epochs and is still decreasing.

\begin{table}[!ht]
\caption{Circuit stability under a fixed rule and reorganization after switching ($p=31$, three seeds; standard deviations in parentheses). Fixed-rule runs show participation after $20{,}000$ and $200{,}000$ epochs. Switched runs use $20{,}000$ epochs in the first phase and $30{,}000$ in the second and show participation after each phase; the reverse switch has not stabilized by the end. Same-rule continuations control for the additional training.}
\label{tab:long}
\begin{center}
\small
\begin{tabular}{lccc}
\toprule
Condition & Circuit PR before & Circuit PR after & Test \\
\midrule
\multicolumn{4}{l}{\emph{Rule held, $20{,}000$ to $200{,}000$ epochs}} \\
Reference ($\beta{=}0$) & $0.60$ & $0.605\,(0.012)$ & $0.997$ \\
Credit equalization ($\beta{=}1$) & $0.99$ & $0.990\,(0.003)$ & $0.975$ \\
\midrule
\multicolumn{4}{l}{\emph{Rule switched after phase one}} \\
$\beta{=}0 \to \beta{=}1$ & $0.613\,(0.016)$ & $0.991\,(0.000)$ & $0.997$ \\
$\beta{=}1 \to \beta{=}0$ & $0.990\,(0.003)$ & $0.722\,(0.016)$ & $0.990$ \\
$\beta{=}0 \to \beta{=}0$ (control) & $0.613\,(0.016)$ & $0.610\,(0.012)$ & $0.997$ \\
$\beta{=}1 \to \beta{=}1$ (control) & $0.990\,(0.003)$ & $0.990\,(0.003)$ & $0.975$ \\
\bottomrule
\end{tabular}
\end{center}
\end{table}

\subsection{When the circuit concentrates}
\label{app:timing}
We train the reference (modular addition, $p = 31$, five fresh seeds) for $20{,}000$ epochs and record circuit participation every $50$ epochs. All runs fit the training set by epoch $100$, when participation is $0.997$ or higher. Test accuracy first reaches $0.95$ between epochs $1250$ and $1500$, and participation has fallen to between $0.59$ and $0.67$ by then. From the grokking epoch to the end of training it changes by $-0.04$ to $+0.05$, while the training loss settles at about $0.04$. The concentration of the circuit thus forms during the grokking delay, together with the generalizing computation, and not in a later phase of training.

\subsection{Gradient and weight decay under Adam}
\label{app:decay}
Our reference adds the coupled decay term $\lambda \vw_j$ to each neuron's incoming-weight gradient $\vg_j$ before Adam, whose update is nearly invariant to a constant rescaling of its input. Full credit equalization multiplies $\vg_j$ by $s_j = \bar{g}/\|\vg_j\|_2$ and leaves the decay term unchanged, so it also changes the ratio between the two. We separate the two effects with four conditions that differ only in how each first-layer neuron's gradient and decay term are scaled before Adam, with $s_j$ recomputed at every step (Table~\ref{tab:decay}). The decay is added after clipping, where Adam adds coupled decay, and biases and the other layers keep the decay $\lambda$. Scaling only the decay term, by $1/s_j$, gives the participation of full equalization, $0.993$, no nearly unused neurons and nearly the same deletion tolerance, and it is more distributed still on the activation-weighted and ablation measures ($0.981$ and $0.924$ against $0.850$ and $0.832$). Scaling both the gradient and the decay term by $s_j$ keeps each neuron's ratio of gradient to decay as in the reference and gives the reference circuit, with participation $0.574$ and $92$ nearly unused neurons. Under Adam with coupled weight decay, the intervention therefore acts through this ratio. Adam's invariance holds exactly only for a constant factor, and $s_j$ changes from step to step, so the two conditions do not give identical updates; the controls measure their effect on the circuit directly.

\paragraph{Where the three non-gradient terms enter the update.} For the incoming row $\vw_j$ of neuron $j$ with loss gradient $\vg_j$ and learning rate $\eta$, Adam with coupled decay forms its moment estimates $\hat{\vm}_j$ and $\hat{\vv}_j$ from $\vg_j + \lambda \vw_j$ and steps
\begin{equation}
\Delta \vw_j = -\eta\, \hat{\vm}_j \big/ \big(\sqrt{\hat{\vv}_j} + \epsilon\big).
\end{equation}
Adam removes a constant overall scale of its input, so a neuron whose gradient is small relative to $\lambda \vw_j$ still receives a step of full size, and that step mostly shrinks its weights. AdamW forms the moments from $\vg_j$ alone and applies the decay outside,
\begin{equation}
\Delta \vw_j = -\eta\, \big[\hat{\vm}_j \big/ \big(\sqrt{\hat{\vv}_j} + \epsilon\big) + \lambda \vw_j\big].
\end{equation}
While $\sqrt{\hat{\vv}_j} \gg \epsilon$ the first term does not depend on the scale of the gradient. Once $\sqrt{\hat{\vv}_j} \ll \epsilon$, as in the wide transformer of Appendix~\ref{app:transformer}, it is approximately $\hat{\vm}_j / \epsilon$, so the gradient competes with $\epsilon \lambda \vw_j$, and $\epsilon$ sets the scale against which the gradient is compared, as $\lambda$ does under coupled decay. A trainable gain $\gamma_j$ and the outgoing column $\va_j$ are updated by the first rule with their own gradients, so each has its own ratio of gradient to decay, which rescaling $\vg_j$ does not change (Appendix~\ref{app:normlayers}). In each case the controls of Table~\ref{tab:pathways} multiply the term that does not scale with the gradient, $\lambda$, $\epsilon$ or the decay of $\gamma_j$ and $\va_j$, by $1/s_j$, or multiply the gradient and this term together by $s_j$.

\begin{table}[!ht]
\caption{Gradient and coupled decay under Adam (modular addition, $p=31$, five seeds, $20{,}000$ epochs; standard deviations in parentheses). $s_j$ is the factor of full credit equalization for neuron $j$, recomputed at every step; the first two rows are the reference and full credit equalization. Nearly unused neurons have an outgoing norm below a tenth of the layer mean. Act.\ and abl.\ are the activation-weighted and ablation participation of Appendix~\ref{app:functional}; random and norm-ranked are deletion tolerances (Appendix~\ref{app:deletion}).}
\label{tab:decay}
\begin{center}
\footnotesize
\setlength{\tabcolsep}{3pt}
\begin{tabular}{llccccccc}
\toprule
Gradient & Decay & Test & Circuit PR & Nearly unused & Act. & Abl. & Random & Norm-ranked \\
\midrule
$\vg_j$ & $\lambda \vw_j$ & $0.995$ & $0.606\,(0.019)$ & $83.6$ & $0.355$ & $0.429$ & $0.675$ & $0.340$ \\
$s_j \vg_j$ & $\lambda \vw_j$ & $0.976$ & $0.993\,(0.001)$ & $0$ & $0.850$ & $0.832$ & $0.776$ & $0.585$ \\
$\vg_j$ & $\lambda \vw_j / s_j$ & $0.979$ & $0.993\,(0.001)$ & $0$ & $0.981$ & $0.924$ & $0.782$ & $0.611$ \\
$s_j \vg_j$ & $s_j \lambda \vw_j$ & $0.996$ & $0.574\,(0.046)$ & $92.4$ & $0.283$ & $0.378$ & $0.664$ & $0.316$ \\
\bottomrule
\end{tabular}
\end{center}
\end{table}
